% ============ 04. Sistemas nacionales de evaluación ============
\section{Sistemas nacionales de información y evaluación: el modelo colombiano en perspectiva}

\subsection{Qué dice la literatura}

La literatura en español y portugués documenta que, entre los países
latinoamericanos, Brasil y Colombia destacan por sus \textbf{sistemas de
información curricular consolidados} — \textbf{Lattes/Plataforma Sucupira} en
Brasil y \textbf{ScienTI-CVLAC/Publindex} en Colombia —, si bien existen otros
registros nacionales (CVar en Argentina, CVU en México); la afirmación de que
son los únicos dos "consolidados" proviene de una fuente comparativa ya clásica
(Revista CTS, ~2010; sección 13.2) y debe leerse con esa fecha. La herramienta
brasileña \emph{ScriptLattes} \citep{menachalco2009scriptlattes} es un
referente metodológico del aprovechamiento analítico de estos registros, y el
\textbf{método del currículum vitae (CV method)} de \citet{canibano2009curriculum}
formaliza su uso en estudios de política científica. Sobre Colombia existen
análisis críticos recientes del \textbf{modelo de categorización de grupos e
investigadores ante el SNCTI} (convocatorias de MinCiencias; estudio en
\emph{Universidad y Sociedad}, sección 13.2), que señalan tensiones entre
clasificación y calidad, sesgos disciplinares y efectos sobre la carrera
académica; y estudios sobre ciencia abierta y métricas de nueva generación en
el contexto de Publindex y ScienTI (sección 13.2), enmarcados en la política
nacional de ciencia abierta (CONPES 4069 de 2021). A nivel internacional, el
diseño de sistemas de evaluación nacional está siendo reformado bajo los
principios de Leiden/CoARA \citep{coara2022} y la literatura de
\emph{research evaluation systems} compara modelos (indicadores vs.\ peer
review vs.\ enfoques mixtos).

\subsection{Relación con este repositorio}

El repositorio audita precisamente el \emph{producto registral} del modelo
colombiano: el padrón de investigadores reconocidos por convocatoria. Sus
hallazgos (cobertura selectiva de variables de diversidad, etiquetado erróneo
de convocatorias, atípicos de edad, captura cambiante de co-filiación, 3 de 6
convocatorias en el consolidado versionado) constituyen evidencia empírica
sobre la \textbf{calidad del registro}, un tema central en la literatura
regional. La comparación implícita del proyecto con las convocatorias como
"ventanas" es coherente con cómo los estudios de la región tratan a Lattes y
ScienTI.

\subsection{Mejoras recomendadas}

\begin{enumerate}
  \item Incorporar explícitamente la \textbf{comparación Brasil--Colombia}
        (Lattes vs.\ ScienTI/CVLAC) como marco interpretativo: qué variables
        captura cada sistema, cómo manejan la desambiguación de autores y la
        actualización curricular, y qué lecciones aplican al diseño de la
        captura de datos de MinCiencias.
  \item Citar y contrastar los \textbf{estudios críticos del modelo colombiano}
        (sección 13.2) para posicionar los hallazgos del observatorio en el
        debate de política, no solo como ejercicio descriptivo.
  \item Alinear las recomendaciones con la agenda de \textbf{métricas
        responsables y ciencia abierta} del país (Publindex, CONPES de ciencia
        abierta), de modo que el observatorio sirva como insumo de política
        pública reproducible.
\end{enumerate}
